% from ega3-1-fr.tex
\oldpage[III]{113}
Ceci va nous permettre de démontrer
(\hyperref[III.2.6.2-fr]{2.6.2}) par récurrence noethérienne
(0\textsubscript{I}, 2.2.2) dans le cas où la restriction de $g$ à
$g^{-1}(X-Z)$ est un isomorphisme sur $X-Z$ (ce qui est suffisant pour
notre propos, comme on l'a vu en
(\hyperref[III.2.6.2.1-fr]{2.6.2.1}))~: il suffira d'établir que si, pour
tout sous-préschéma fermé $X_1$ de $X$, dont l'espace sous-jacent est
$\neq X$, la conclusion de (\hyperref[III.2.6.2-fr]{2.6.2}) est vraie
pour le faisceau $\mathcal L_1$, alors elle est vraie aussi pour le
faisceau $\mathcal L$.

\begin{env}[2.6.2.4]
\label{III.2.6.2.4-fr}
Soient alors $\mathcal A=\mathcal O_X$, $\mathcal B=g_*(\mathcal O_{X'})$,
de sorte que $\mathcal B$ est une sous-$\mathcal A$-Algèbre de
$\mathcal R(X)$, qui est un $\mathcal A$-Module cohérent~; en outre, la
restriction $\mathcal B|(X-Z)$ est égale à $\mathcal A|(X-Z)$. Soit
$\mathcal K$ le conducteur de $\mathcal B$ sur $\mathcal A$, c'est-à-dire
le plus grand sous-$\mathcal A$-Module quasi-cohérent de $\mathcal A$ tel
que $\mathcal B\mathcal K\subset\mathcal A$ (ou encore l'annulateur du
$\mathcal A$-Module $\mathcal B/\mathcal A$
(0\textsubscript{I}, 5.3.7)), ce qui entraîne
$\mathcal B\mathcal K=\mathcal K$. Il est clair que
$\mathcal K_x=\mathcal A_x$ en tous les points admettant un voisinage
$W_x$ tel que $g$ soit un isomorphisme de $g^{-1}(W_x)$ sur $W_x$, et en
particulier en tous les points de $X-Z$ et dans un voisinage de tout
point générique d'une composante irréductible de $X$. Considérons alors
le sous-préschéma fermé $Z_1=\operatorname{Spec}(\mathcal A/\mathcal K)$
de $X$ défini par $\mathcal K$~; il est encore propre sur $Y$ car le
sous-espace $Z_1$ est fermé dans $Z$ (II, 5.4.10). De plus, la définition
de $\mathcal K$ montre que
$\mathcal B|(X-Z_1)=\mathcal A|(X-Z_1)$~; on voit donc qu'on peut toujours
se ramener au cas où
$Z=\operatorname{Spec}(\mathcal A/\mathcal K)$, et comme on a vu que
$X-Z_1$ est un ouvert non vide de $X$, on peut toujours supposer que
l'espace $Z$ est distinct de $X$.
\end{env}

\begin{env}[2.6.2.5]
\label{III.2.6.2.5-fr}
Considérons $X'$ comme égal à $\operatorname{Spec}(\mathcal B)$ (puisque
$g$ est affine) et soit $\mathcal K'=\widetilde{\mathcal K}$, Idéal
cohérent de $\mathcal O_{X'}$ tel que $g_*(\mathcal K')=\mathcal K$
(II, 1.4.1)~; le sous-préschéma fermé
$Z'=g^{-1}(Z)=Z\times_X X'$ de $X'$ est défini par $\mathcal K'$ et égal
à $\operatorname{Spec}(\mathcal B/\mathcal K)$ (II, 1.4.10)~; comme
$h:Z'\to Z$ est un morphisme fini (II, 6.1.5, (iii)), $Z'$ est propre
sur $Y$ (II, 6.1.11 et 5.4.2, (ii)).

Cela posé, il nous faut prouver que pour tout $x\in X$ et tout voisinage
ouvert $U$ de $x$, il existe une section $s$ d'un
$\mathcal L^{\otimes n}$ $(n>0)$ au-dessus de $X$ telle que
$x\in X_s\subset U$ (II, 4.5.2)~; nous distinguerons deux cas~:

$1^\circ$ On a $x\in X-Z$~; on peut évidemment supposer alors que l'on a
aussi $U\subset X-Z$, donc l'ouvert $U'=g^{-1}(U)$ ne rencontre pas
$Z'$. Comme $\mathcal L'$ est ample par hypothèse, il existe un $n>0$ et
une section $s'$ de $\mathcal L'^{\otimes n}$ au-dessus de $X'$ telle
que $x'=g^{-1}(x)\in X'_{s'}\subset g^{-1}(U)$ (II, 4.5.2). En outre, on
peut supposer que $\mathcal K'\otimes\mathcal L'^{\otimes n}$ soit
engendré par ses sections au-dessus de $X'$ (II, 4.5.5), donc, comme
$\mathcal K'_{x'}=\mathcal O_{x'}$, il y a une de ces sections $s''$
telle que $s''(x')\neq0$~; en la multipliant par $s'$ (ce qui revient à
remplacer $n$ par $2n$), on voit qu'on peut supposer aussi que
$x'\in X'_{s''}\subset g^{-1}(U)$. Cela étant, il résulte de
(0\textsubscript{I}, 5.4.10) que l'on a un isomorphisme canonique
\[
\Gamma(X,\mathcal K\otimes\mathcal L^{\otimes n})
\xrightarrow{\sim}
\Gamma(X',\mathcal K'\otimes\mathcal L'^{\otimes n}).
\]
La section $s$ de $\mathcal K\otimes\mathcal L^{\otimes n}$ qui
correspond à $s''$ par cet isomorphisme a évidemment les propriétés
voulues.

$2^\circ$ On a $x\in Z$. Soit $\mathcal J$ l'Idéal cohérent de
$\mathcal O_X$ définissant le sous-préschéma fermé réduit de $X$ ayant
pour espace sous-jacent $X-U$, et considérons dans $\mathcal B$ les deux

